% ===========================================================================
\chapter{A place to work}\label{ch:workspace}

The agent works in one folder, the \emph{workspace}, and nowhere else.
Every path the model gives is taken to be relative to it, and a path that
leads outside it is refused. This is the first of the agent's safety rules,
and the most important: a model can make mistakes, and it can be misled
(a file it reads might contain instructions written to trick it), so what
it can reach should be limited by the program, not by the model's good
sense.

\section{Paths that stay inside}

The rule sounds simple: refuse a path that starts with \code{/}, or that
contains \code{..}. But \code{..} is not always a problem
(\code{src/../loans.nex} is just \code{loans.nex}), and other paths escape
without looking as if they do (\code{notes/../../x}). The reliable test is
to do what the operating system will do: put the path after the workspace
folder, resolve every \code{.} and \code{..} in it, and see whether the
result is still inside. \code{Path.normalize} does the resolving.

\begin{newfile}{workspace.nex}
-- workspace.nex: the folder the agent works in, and the rule that it
-- never reaches outside it.

intern io/Path

class Workspace
create
  -- The workspace rooted at `folder`, which must exist.
  make(folder: String)
  require
    is_a_folder: (create Path.make(folder)).is_directory()
  do
    root := (create Path.make(folder)).absolute().normalize().path
  end

feature
  root: String    -- the workspace folder, as an absolute path

  -- Whether `relative`, a path the model gave, names something inside
  -- the workspace. "loans.nex", "src/a.nex" and "." do; "/etc/passwd"
  -- and "../secrets" do not, and neither does "a/../../b", which only
  -- looks as if it stays inside.
  contains(relative: String): Boolean do
    if relative = "" or relative.starts_with("/") or relative.starts_with("~") then
      result := false
    else
      result := is_under_root(joined(relative))
    end
  end
\end{newfile}

A path starting with \code{~} is refused too: bash would read it as your
home folder.

Every other routine of \code{Workspace} takes a path the model gave, and
every one of them has the same precondition: \code{inside:
contains(relative)}. That is the rule, stated once per routine, and
checked on every call. A tool that forgets to check a path first doesn't
quietly read \code{/etc/passwd}; it stops with \code{Precondition
violation: inside}. In the next chapter the tools will check paths
themselves, to give the model a helpful message, but the contract is what
guarantees it.

\begin{continued}{workspace.nex}
  -- The absolute path of `relative`.
  full_path(relative: String): String
  require
    inside: contains(relative)
  do
    result := joined(relative)
  ensure
    under_root: is_under_root(result)
  end

  -- The path of `full`, a path inside the workspace, relative to the
  -- root: "loans.nex", or "." for the root itself.
  relative_path(full: String): String
  require
    under_root: is_under_root(full)
  do
    if full = root then
      result := "."
    else
      result := full.substring(root.length + 1, full.length)
    end
  end

  exists(relative: String): Boolean
  require
    inside: contains(relative)
  do
    result := (create Path.make(full_path(relative))).exists()
  end

  is_file(relative: String): Boolean
  require
    inside: contains(relative)
  do
    result := (create Path.make(full_path(relative))).is_file()
  end

  is_folder(relative: String): Boolean
  require
    inside: contains(relative)
  do
    result := (create Path.make(full_path(relative))).is_directory()
  end
\end{continued}

\section{Reading and writing}

Reading needs the path to be a file. Writing creates any folders on the
way, so that the model can write \code{notes/plan.md} without first
making \code{notes}. Its postcondition reads the file back: what was
written is what is there.

\begin{continued}{workspace.nex}
  read(relative: String): String
  require
    inside: contains(relative)
    a_file: is_file(relative)
  do
    result := (create Path.make(full_path(relative))).read_text()
  end

  -- Write `text` to the file `relative`, replacing it if it exists and
  -- creating any folders on the way to it.
  write(relative, text: String)
  require
    inside: contains(relative)
    not_a_folder: not is_folder(relative)
  do
    let file: Path := create Path.make(full_path(relative))
    let parent: ?Path := file.parent()
    if ?parent as folder then
      folder.create_directories()
    end
    file.write_text(text)
  ensure
    written: read(relative) = text
  end
\end{continued}

\code{file.parent()} has the type \code{?Path}: a \emph{detachable}
\code{Path}, which may be \code{nil} (the root of the file system has no
parent). Nex won't let us call anything on it until we have checked.
\code{if ?parent as folder then} does both at once: the branch runs only
if \code{parent} is not \code{nil}, and inside it \code{folder} is the
same value as a plain \code{Path}.

\section{What is there}

To find its way around a project, the model needs to see what files it
has. \code{entries} lists everything below a folder, folders first marked
with a \code{/}; \code{files} only the files, which is what a search will
want. Both leave out hidden files and folders, whose names start with a
dot: a \code{.git} folder holds thousands of files that are of no use to
the model and would bury the ones that are.

\begin{continued}{workspace.nex}
  -- Everything below the folder `relative`, as paths relative to the
  -- root, in alphabetical order. Folders end in "/". Hidden files and
  -- folders (".git", ".cache", ...) are left out.
  entries(relative: String): Array[String]
  require
    inside: contains(relative)
    a_folder: is_folder(relative)
  do
    collect(create Path.make(full_path(relative)), result, true)
  end

  -- The files below the folder `relative`, like entries but without the
  -- folders themselves.
  files(relative: String): Array[String]
  require
    inside: contains(relative)
    a_folder: is_folder(relative)
  do
    collect(create Path.make(full_path(relative)), result, false)
  end
\end{continued}

The private routines do the work. \code{collect} walks the folder tree,
depth first. The operating system lists a folder's contents in no
particular order, so \code{collect} sorts them by name first: the same
folder always gives the same list, which makes the model's view of a
project, and our checks, repeatable. \code{sort} takes a function that
compares two elements, here written in place with \code{fn}.

\begin{continued}{workspace.nex}
  private feature
    -- `relative` joined to the root, with "." and ".." resolved.
    joined(relative: String): String do
      result := (create Path.make(root + "/" + relative)).normalize().path
    end

    is_under_root(full: String): Boolean do
      result := full = root or full.starts_with(root + "/")
    end

    -- Add what is below `folder` to `found`, depth first, in order.
    collect(folder: Path, found: Array[String], with_folders: Boolean) do
      let children: Array[Path] := folder.list().sort(fn (a, b: Path): Integer do
        result := a.name().compare(b.name())
      end)
      across children as child do
        if not child.name().starts_with(".") then
          if child.is_directory() then
            if with_folders then
              found.add(relative_path(child.normalize().path) + "/")
            end
            collect(child, found, with_folders)
          else
            found.add(relative_path(child.normalize().path))
          end
        end
      end
    end

invariant
  root_is_absolute: (create Path.make(root)).absolute().normalize().path = root
end
\end{continued}

Why \code{root + "/"} in \code{is_under_root}, and not just \code{root}?
Because the folder \code{/home/ada/agent-old} starts with
\code{/home/ada/agent} too.

\section{A fresh copy}

The agent will change the files it works on. To be able to run it again
and again on the same starting point, it works on a copy of the sample
project, made fresh each time:

\begin{append}{workspace.nex}
-- A workspace holding a fresh copy of the folder `template`, at `folder`.
-- Whatever was at `folder` before is deleted first, so every run starts
-- from the same files.
function fresh_workspace(template, folder: String): Workspace
require
  template_exists: (create Path.make(template)).is_directory()
  not_the_template: (create Path.make(template)).absolute().normalize().path
    /= (create Path.make(folder)).absolute().normalize().path
do
  let target: Path := create Path.make(folder)
  if target.exists() then
    target.delete_tree()
  end
  (create Path.make(template)).copy_to(target)
  result := create Workspace.make(folder)
end
\end{append}

The second precondition guards against the one mistake here that would
hurt: deleting the template itself.

\begin{tryit}{try_workspace.nex}
intern Workspace
intern Helpers

let ws: Workspace := fresh_workspace(user_path("sample_project"), user_path("scratch"))
across ["loans.nex", ".", "src/../loans.nex", "../secrets", "/etc/passwd",
    "notes/../../x", "~/.ssh"] as path do
  console.print_line(path.pad_end(" ", 20) + ws.contains(path).to_string)
end
ws.write("notes/plan.md", "1. run the checks\n")
console.print_line(ws.entries(".").to_string)
console.print_line(ws.read("notes/plan.md"))

let refused: String := ""
do
  console.print_line(ws.read("../helpers.nex"))
rescue
  refused := "" + exception
end
console.print_line(refused)
\end{tryit}
\tryoutput{try_workspace}

\begin{aside}
\textbf{What this does not stop.} A path is checked as text. If the
workspace contained a symbolic link to a folder outside it, a path through
the link would pass the test. Our sample project has none, and a more
careful agent would resolve links too (Java's \code{toRealPath} does). And
the rule covers the file tools only: a shell command can still
\code{cat /etc/passwd}. That is why running commands will need permission
(Chapter~\ref{ch:permission}).
\end{aside}
